\documentclass[12pt,a4paper]{article}
\usepackage{float}

%% ── Packages ──────────────────────────────────────────────────
\usepackage[utf8]{inputenc}
\usepackage[T1]{fontenc}
\usepackage[british]{babel}
\usepackage{amsmath,amssymb,amsthm}
\usepackage{mathtools}
\usepackage{geometry}
\usepackage{hyperref}
\usepackage{xcolor}
\usepackage{graphicx}
\usepackage{booktabs}
\usepackage{hyphenat}
\usepackage{array}
\usepackage{enumitem}
\usepackage{microtype}
\usepackage[numbers]{natbib}
\usepackage{tcolorbox}
\usepackage{tikz}
\usetikzlibrary{arrows.meta,positioning,shapes.geometric}

\geometry{margin=2.5cm}
\hypersetup{colorlinks=true,linkcolor=blue!70!black,
            citecolor=green!50!black,urlcolor=blue!70!black}

%% ── tcolorbox environments ────────────────────────────────────
\tcbuselibrary{theorems,skins,breakable}

\newtcbtheorem[number within=section]{theorem}{Theorem}%
{colback=blue!5,colframe=blue!40!black,fonttitle=\bfseries,
 breakable}{thm}

\newtcbtheorem[number within=section]{conjecture}{Conjecture}%
{colback=orange!5,colframe=orange!60!black,fonttitle=\bfseries,
 breakable}{conj}

\newtcbtheorem[number within=section]{prediction}{Prediction}%
{colback=green!5,colframe=green!50!black,fonttitle=\bfseries,
 breakable}{pred}

\newtcbtheorem[number within=section]{openproblem}{Open Problem}%
{colback=red!5,colframe=red!50!black,fonttitle=\bfseries,
 breakable}{op}

\newtcbtheorem[number within=section]{definition}{Definition}%
{colback=gray!5,colframe=gray!50,fonttitle=\bfseries,
 breakable}{def}

%% ── Epistemic status box ──────────────────────────────────────
\newenvironment{epistemicbox}{%
  \begin{tcolorbox}[colback=yellow!8,colframe=yellow!60!black,
    title={\bfseries Epistemic Status},breakable]}%
  {\end{tcolorbox}}

%% ── Macros ────────────────────────────────────────────────────
\newcommand{\epsgeom}{\varepsilon_{\mathrm{geom}}}
\newcommand{\DEN}{\Delta\mathrm{EN}}
\newcommand{\Gloc}{G_{\mathrm{loc}}}
\newcommand{\Gglob}{G_{\mathrm{glob}}}
\newcommand{\Hloc}{H_{\mathrm{loc}}}
\newcommand{\ZZ}{\mathbb{Z}}
\newcommand{\QQ}{\mathbb{Q}}

%% ══════════════════════════════════════════════════════════════
\begin{document}

\title{\textbf{D-exp-SP2: Topological Optimality Criterion\\
for Coherent Photon-to-Electron Conversion}\\[6pt]
{\large Projective Dynamic Logo Framework — Exploratory Document}}

\author{C\'edric Laubscher\\[4pt]
\small Independent Researcher, Switzerland\\
\small \href{https://cedriclaubscher.ch}{cedriclaubscher.ch}\\
\small \href{mailto:cedric@cedriclaubscher.ch}{cedric@cedriclaubscher.ch}}

\date{May 2026\\[4pt]
\small PDL corpus version: D01--D53 + D-exp-SP2\\
\small Zenodo community: \texttt{pdl-framework}}

\maketitle

%% ── Epistemic status ──────────────────────────────────────────
\begin{epistemicbox}
This is an \textbf{exploratory document}. It is not part of the PDL core
corpus (D01--D53). Its results are classified as follows:
\begin{itemize}[noitemsep]
  \item \textbf{Theorems}: Theorem~\ref{thm:pdlsp2} (PDL-SP2) and its
    corollaries — proved analytically and verified exhaustively.
  \item \textbf{Observations}: Rule~A ($k_{\mathrm{avg}}=4$,
    $\DEN>1.0\Rightarrow$ direct gap) — verified on 15 materials,
    not a theorem.
  \item \textbf{Conjectures}: PDL-H (Conjecture~\ref{conj:pdlh}) —
    verified exhaustively on $2^{24}$ cases, not yet proved analytically.
  \item \textbf{Predictions}: Falsifiable experimental predictions
    (Section~\ref{sec:predictions}) — no free parameters.
\end{itemize}
No material was presupposed optimal prior to computation.
\end{epistemicbox}

\begin{abstract}
We derive a topological optimality criterion for coherent photon-to-electron
conversion from the axioms C1--C4 of the Projective Dynamic Logo (PDL)
framework. The central result (Theorem PDL-SP2) establishes that
the minimal geometric leakage $\epsgeom = 0$ is achievable on a bond
graph $G$ if and only if $G$ is bipartite — recovering the Harary
balance theorem (1953) as a corollary of C1--C4 without presupposition.
A four-criterion composite score derived from PDL selects six optimal
inorganic candidates from 21 families: BP, AlAs, GaP, AlP, GaN, and SiC.
GaN and BP are identified as priority candidates with direct gaps
of 1.73~eV and 1.24~eV respectively. An extension to hierarchical graphs
(Conjecture PDL-H) shows that global bipartition is the sole determining
condition — local graph structure is irrelevant for $\epsgeom(H)=0$.
This result is verified exhaustively on $2^{24} = 16\,777\,216$ sign
assignments for the case $H(C_5, C_4)$. Biological connection:
the Mg--Mg transfer network of Photosystem~II is bipartite, consistent
with PDL-H, while individual chlorophyll molecules are not — a two-scale
architecture that PDL explains naturally. Three falsifiable experimental
predictions are formulated, testable by femtosecond pump-probe spectroscopy
without free parameters.
\end{abstract}

\clearpage

\tableofcontents
\newpage

%% ══════════════════════════════════════════════════════════════
\section{Introduction and Motivation}
\label{sec:intro}

The conversion of photons to free electrons with near-unity quantum
efficiency is the central challenge of photovoltaic science and the
defining achievement of biological photosynthesis. Photosystem~II achieves
charge separation in under 3~ps with a quantum yield approaching 100\%
\cite{Fleming1988}. Current best photovoltaic materials operate far below
this limit, and no theoretical framework has derived from first principles
which structural features make a material capable of achieving such
efficiency.

The Projective Dynamic Logo (PDL) framework \cite{PDL_DM_v20} derives
fundamental physical structures from four combinatorial axioms (C1--C4)
on finite signed graphs, without presupposing spacetime, particles, or
fields. Its central object is the graph $K_4$, identified as the minimal
admissible closure. Document D49 of the PDL corpus derives the London
equation from C1--C4 via the metric identity
$\nu = \tfrac{1}{4}\|\psi_1 - T^k\psi_1\|^2$, characterising coherent
transport as a norm of the difference between a state and its image under
$K_4$ pulsation. Document D43 defines the geometric leakage parameter
$\epsgeom(G,v)$ as the fraction of cycles of length $\leq 6$ incident
to node $v$ with negative sign product.

This document applies these PDL tools to a concrete problem: identifying,
from first principles, which materials and material architectures are
topologically optimal for coherent photon-to-electron conversion.
The approach is strictly deductive — no material is presupposed optimal,
and all candidate selection follows from the axioms.

%% ══════════════════════════════════════════════════════════════
\section{PDL Ingredients}
\label{sec:ingredients}

We recall the four PDL results used throughout this document.
All are theorems of the core corpus; their proofs are not reproduced here.

\begin{definition}{Geometric Leakage (D43)}{eps_geom_def}
Let $G = (V, E, s)$ be a finite signed graph with $s: E \to \{+1,-1\}$.
For a node $v \in V$, define
\[
  \epsgeom(G, s, v)
  = \frac{\bigl|\{C \in \mathcal{C}_{\leq 6}(v) :
    \prod_{e \in C} s(e) = -1\}\bigr|}
         {|\mathcal{C}_{\leq 6}(v)|}
\]
where $\mathcal{C}_{\leq 6}(v)$ denotes the set of simple cycles of
length $\leq 6$ incident to $v$. The global leakage is
$\epsgeom(G,s) = \frac{1}{|V|}\sum_{v} \epsgeom(G,s,v)$.
\end{definition}

\begin{definition}{Bipartite Graph}{bip_def}
A graph $G = (V, E)$ is \emph{bipartite} if $V$ can be partitioned
into two disjoint sets $A$ and $B$ such that every edge connects a vertex
in $A$ to a vertex in $B$. Equivalently, $G$ contains no cycle of odd
length.
\end{definition}

The four PDL results used here are:
\begin{enumerate}[label=\textbf{P\arabic*.}]
  \item \textbf{Minimal closure (D16a)}: $K_4$ is the smallest admissible
    graph satisfying C1--C4.
  \item \textbf{Geometric leakage (D43)}: $\epsgeom(p) = 329/10087$
    and $\epsgeom(n) = 468/9960$ are unconditional theorems of C1--C4.
  \item \textbf{Coherent transport (D49)}: The London equation is derived
    from C1--C4 via $\nu = \tfrac{1}{4}\|\psi_1 - T^k\psi_1\|^2$.
  \item \textbf{Connectivity (C1)}: Every PDL-admissible graph is
    finite and connected. This forces the number of connected components
    $c = 1$, and hence $|\mathrm{eps}_0| = 2^{N-1}$ (see below).
\end{enumerate}

%% ══════════════════════════════════════════════════════════════
\section{Theorem PDL-SP2: Bipartition and Zero Leakage}
\label{sec:theorem}

\begin{theorem}{PDL-SP2 — Bipartition and Zero Leakage}{pdlsp2}
Let $G$ be a connected graph with $N$ nodes and $c = 1$ connected
component. A sign assignment $s: E \to \{+1,-1\}$ satisfies
$\epsgeom(G,s) = 0$ if and only if there exists a colouring
$f: V \to \{+1,-1\}$ such that $s(u,v) = f(u)\cdot f(v)$
for every edge $(u,v) \in E$.

Equivalently: $\epsgeom(G,s) = 0$ is achievable if and only if
$G$ is bipartite.
\end{theorem}

\begin{proof}[Proof sketch]
\textbf{Part A} (every $s_f$ satisfies $\epsgeom = 0$): For any cycle
$(v_1,\ldots,v_k,v_1)$ and any colouring $f$,
\[
  \prod_{i=1}^{k} s_f(v_i,v_{i+1})
  = \prod_{i=1}^{k} f(v_i)\cdot f(v_{i+1})
  = \prod_{i=1}^{k} f(v_i)^2 = +1,
\]
since each $f(v_i) \in \{+1,-1\}$ appears exactly twice. Hence every
cycle has positive sign product and $\epsgeom = 0$.

\textbf{Part B} (reconstruction): Conversely, given $s$ with
$\epsgeom = 0$, fix a root node $r$ with $f(r) = +1$. For any
other node $v$, choose a path $P(r \to v)$ and define
$f(v) = \prod_{e \in P} s(e)$. Since every cycle has product $+1$,
this definition is path-independent (connectivity is essential here).
One verifies $s(u,v) = f(u)\cdot f(v)$ for every edge.
\end{proof}

\textbf{Connection to Harary (1953)}: Theorem~\ref{thm:pdlsp2} recovers
the Harary balance theorem \cite{Harary1953} as a consequence of C1--C4,
without presupposing it.

\begin{theorem}{Counting Optimal Assignments}{counting}
For a connected graph $G$ with $N$ nodes, $M$ edges, and $k$-regular
with $k=3$ (3-regular, $M = 3N/2$):
\[
  |\{\text{sign assignments with } \epsgeom = 0\}|
  = 2^{N-1}, \qquad
  \text{fraction} = \frac{2^{N-1}}{2^{3N/2}} = 2^{-N/2-1}.
\]
The bijection $f \mapsto s_f$ is 2-to-1 (since $f$ and $-f$ give
the same assignment), accounting for the factor of $1/2$.
\end{theorem}

\textbf{Exhaustive verification}: Both theorems were verified
exhaustively for all 3-regular graphs with $N = 4, 6, 8, 10$ nodes
(7 distinct graphs, all sign assignments enumerated), confirming
zero exceptions. The reconstruction algorithm (Part~B) was verified on
all optimal assignments for each graph, with zero failures.

For non-connected graphs with $c$ components, the counting formula
generalises to $2^{N-c}$, verified on the disjoint union of two $K_4$
graphs ($N=8$, $c=2$, $2^6 = 64$ optimal assignments confirmed).

%% ══════════════════════════════════════════════════════════════
\section{Optimality Criterion for Photon-to-Electron Conversion}
\label{sec:criterion}

Theorem~\ref{thm:pdlsp2} establishes bipartition as the necessary and
sufficient condition for zero geometric leakage. For application to
photovoltaic materials, three additional conditions are derived from
physical considerations anchored in PDL.

\begin{definition}{PDL-SP2 Composite Score}{score_def}
A material with bond graph $G$, band gap $E_g$, and bipartition
groups $A$, $B$ receives a score $\sigma \in \{0,1,2,3,4\}$:
\begin{align*}
  \sigma &= \mathbf{1}[G\text{ bipartite}]
           + \mathbf{1}[G\text{ connected}]
           + \mathbf{1}[1.1 \leq E_g \leq 2.5\text{ eV}]
           + \mathbf{1}[\text{elements}(A) \\ &\neq \text{elements}(B)].
\end{align*}
Criterion 1 is the PDL theorem; Criterion 3 targets the solar spectrum;
Criterion 4 requires chemical heterogeneity (symmetry breaking) to
create a preferred direction of charge transfer.
\end{definition}

\textbf{Rule A} (Observation, not theorem): Among 15 materials tested,
the rule $k_{\mathrm{avg}} = 4$ and $\DEN > 1.0 \Rightarrow$ direct gap
holds for 4/4 materials satisfying both conditions (GaN, AlN, InN, ZnO).
The sole exception at $\DEN > 1.0$ is SnO$_2$ ($k_{\mathrm{avg}} = 2.67$,
rutile structure), correctly excluded by the $k_{\mathrm{avg}}$ condition.
This observation is not a PDL theorem; DFT validation is required.

%% ══════════════════════════════════════════════════════════════
\section{Material Screening: 21 Candidates from Four Families}
\label{sec:screening}

Bond graphs were constructed from crystallographic structures retrieved
from the Materials Project \cite{Jain2013} using covalent-radius cutoffs
(Alvarez 2008 \cite{Alvarez2008}). Each structure was retrieved dynamically
by chemical formula (most stable phase, $E_{\mathrm{hull}} = 0$), with
systematic element verification to exclude database artefacts.

\subsection{Families Tested}

Four families were selected from PDL optimality conditions, without
presupposing any candidate:
\begin{itemize}[noitemsep]
  \item \textbf{F1}: III-V nitrides and carbides (GaN, AlN, InN, BN,
    GaP, InP, AlP, SiC)
  \item \textbf{F2}: Chalcopyrites and sulfo-salts (CuInSe$_2$,
    CuGaSe$_2$, CuInS$_2$, AgInSe$_2$, ZnSnP$_2$, CuGaS$_2$)
  \item \textbf{F3}: Graphene analogues (BN, BP, AlAs, InAs)
  \item \textbf{F4}: Ternary oxides and spinels (ZnO, CuAlO$_2$,
    CuGaO$_2$, ZnGa$_2$O$_4$, MgAl$_2$O$_4$, Cu$_2$O, SnO$_2$)
\end{itemize}

\subsection{Results}

\begin{table}[h]
\centering
\caption{PDL-SP2 screening results. Score 4/4 = optimal.
Gap type: D = direct, I = indirect (experimental literature).
Bipartition groups listed for score-4 candidates only.}
\label{tab:screening}
\begin{tabular}{lrrllll}
\toprule
Material & Score & $E_g$ (eV) & Type & Sym. group & Group A & Group B \\
\midrule
BP       & 4/4 & 1.24 & I & F-43m & \{B\}  & \{P\}  \\
AlAs     & 4/4 & 1.50 & I & F-43m & \{Al\} & \{As\} \\
GaP      & 4/4 & 1.60 & I & F-43m & \{Ga\} & \{P\}  \\
AlP      & 4/4 & 1.63 & I & F-43m & \{Al\} & \{P\}  \\
GaN      & 4/4 & 1.73 & D & P6$_3$mc & \{Ga\} & \{N\}  \\
SiC      & 4/4 & 1.84 & I & R3m & \{Si\} & \{C\}  \\
\midrule
InP      & 3/4 & 0.45 & D & F-43m & — & — \\
ZnO      & 3/4 & 0.72 & D & P6$_3$mc & — & — \\
AlN      & 3/4 & 4.05 & D & P6$_3$mc & — & — \\
ZnSnP$_2$ & 3/4 & 0.70 & D & I-42d & — & — \\
SnO$_2$  & 3/4 & 0.65 & I & P4$_2$/mnm & — & — \\
\midrule
Cu$_2$O  & 0/4 & 0.51 & — & Pn-3m & \multicolumn{2}{l}{non-bipartite} \\
MgAl$_2$O$_4$ & 0/4 & 5.11 & — & Fd-3m & \multicolumn{2}{l}{non-bipartite} \\
\bottomrule
\end{tabular}
\end{table}

\textbf{Priority candidates}: GaN (score 4/4, direct gap 1.73~eV,
$\DEN = 1.23$, Rule~A satisfied) and BP (score 4/4, gap 1.24~eV
near the Shockley--Queisser optimum of 1.34~eV, largely unexplored
for photovoltaics).

%% ══════════════════════════════════════════════════════════════
\section{Conjecture PDL-H: Hierarchical Graphs}
\label{sec:hierarchical}

Biological photosynthesis operates at two topological scales simultaneously:
individual chlorophyll molecules contain a five-membered ring (cyclopentanone,
girth~$= 5$, non-bipartite), yet the Mg--Mg transfer network between
chlorophylls is bipartite (verified on structure 3ARC of Photosystem~II,
Protein Data Bank \cite{Umena2011}). This two-scale architecture motivates
the following extension of PDL-SP2.

\begin{definition}{Hierarchical Graph}{hier_def}
A hierarchical graph $H = (\Gloc, \Gglob, \phi)$ is defined by:
\begin{itemize}[noitemsep]
  \item $\Gloc$: local graph (intra-chromophore), any structure;
  \item $\Gglob$: global graph (inter-chromophore network);
  \item $\phi \in V(\Gloc)$: connection node (analogue of Mg).
\end{itemize}
The total graph $H$ places one copy of $\Gloc$ on each node of $\Gglob$,
connecting adjacent copies via $\phi$.
\end{definition}

\begin{conjecture}{PDL-H — Global Bipartition Dominates}{pdlh}
For any hierarchical graph $H = (\Gloc, \Gglob, \phi)$ with $\Gloc$
connected:
\[
  \epsgeom(H) = 0 \text{ is achievable}
  \quad \Longleftrightarrow \quad
  \Gglob \text{ is bipartite},
\]
independently of the structure of $\Gloc$ (its girth, bipartition
status, or size).
\end{conjecture}

\textbf{Exhaustive verification}: The case $H(C_5, C_4)$ — pentagonal
local graph (girth~$= 5$, non-bipartite) with square global graph
(bipartite) — was verified exhaustively over all $2^{24} = 16\,777\,216$
sign assignments. Result: $\epsgeom = 0$ is achievable, with
$524\,288 = 2^{19}$ optimal assignments. This is consistent with the
counting formula $2^{N-1}$ (here $N = 20$, $2^{19}$ pairs under the
symmetry $f \sim -f$).

Additional verification: cycles $C_n$ for $n = 3, \ldots, 12$ as local
graphs, combined with three bipartite global graphs (P4, C4, C6).
In all computable cases (exact enumeration or greedy heuristic),
$\epsgeom(H) = 0$ was achieved when $\Gglob$ is bipartite.

\textbf{Biological implication}: Conjecture PDL-H explains why
Photosystem~II operates efficiently despite individual chlorophylls
being non-bipartite. The bipartite Mg--Mg network is sufficient for
$\epsgeom(H) = 0$ at the system level.

\textbf{Design implication}: For artificial photovoltaic systems,
PDL-H implies that the local chromophore structure is topologically
irrelevant — only the global network topology must be bipartite.
This substantially relaxes the material design constraints identified
by PDL-SP2 alone.

%% ══════════════════════════════════════════════════════════════
\section{Falsifiable Predictions}
\label{sec:predictions}

\begin{prediction}{Charge Separation Hierarchy}{P1}
Among the score-4/4 candidates, the charge separation time $\tau_{\mathrm{cs}}$
satisfies the hierarchy
\[
  \tau_{\mathrm{cs}}(\mathrm{GaN})
  < \tau_{\mathrm{cs}}(\mathrm{BP})
  < \tau_{\mathrm{cs}}(\mathrm{SiC})
  < \tau_{\mathrm{cs}}(\mathrm{GaP}),
\]
with $\tau_{\mathrm{cs}}(\mathrm{GaN}) < 3$~ps, testable by femtosecond
transient absorption spectroscopy under identical excitation conditions.
No free parameters. Rationale: GaN satisfies Rule~A (direct gap, $k=4$,
$\DEN > 1.0$); the hierarchy follows from decreasing $\DEN$ and
direct-to-indirect gap transition.
\end{prediction}

\begin{prediction}{Hierarchical Photovoltaic Architecture}{P2}
A van der Waals heterostructure of two score-4/4 materials (e.g.\
GaN/BP or GaN/SiC) with a bipartite inter-layer coupling network will
exhibit charge separation times comparable to Photosystem~II
($\tau_{\mathrm{cs}} < 3$~ps), by Conjecture PDL-H. The external
quantum efficiency will exceed that of either material in isolation.
Testable by EQE measurement and pump-probe spectroscopy.
\end{prediction}

\begin{prediction}{Non-Bipartite Disqualification}{P3}
Materials with non-bipartite bond graphs (Cu$_2$O, MgAl$_2$O$_4$)
will exhibit structurally longer charge separation times than bipartite
candidates under identical geometry and excitation, reflecting the
irreducible leakage $\epsgeom > 0$ imposed by odd cycles in the bond
graph. Testable by comparative femtosecond spectroscopy.
\end{prediction}

%% ══════════════════════════════════════════════════════════════
\section{Epistemic Table}
\label{sec:epistemic}

\begin{table}[h]
\centering
\caption{Epistemic status of all results in this document.}
\label{tab:epistemic}
\begin{tabular}{p{5cm}p{3cm}p{5cm}}
\toprule
Result & Status & Evidence \\
\midrule
Theorem PDL-SP2 ($\epsgeom=0 \Leftrightarrow$ bipartite)
  & Theorem & Analytic proof + exhaustive verification, 7 graphs \\
Counting formula $2^{N-1}$
  & Theorem & Exhaustive, $N=4,6,8,10$ \\
Harary (1953) recovered from C1--C4
  & Theorem & Analytic \\
Score-4/4 candidates (GaN, BP, AlAs, GaP, AlP, SiC)
  & Observation & Materials Project screening, 21 materials \\
Rule~A ($k=4$, $\DEN>1.0 \Rightarrow$ direct gap)
  & Observation & 15 materials, 4/4 correct, 1 exception (SnO$_2$) \\
Conjecture PDL-H
  & Conjecture & Exhaustive $2^{24}$, heuristic $n=3..12$ \\
Predictions P1--P3
  & Prediction & Derived from PDL-SP2 + PDL-H, no free parameters \\
\bottomrule
\end{tabular}
\end{table}

%% ══════════════════════════════════════════════════════════════
\section{Open Problems}
\label{sec:open}

\begin{openproblem}{Analytic Proof of PDL-H}{OP-SP2-1}
Prove Conjecture~\ref{conj:pdlh} analytically for all connected
hierarchical graphs $H = (\Gloc, \Gglob, \phi)$. The exhaustive
verification establishes the result for $H(C_5, C_4)$ and heuristic
evidence covers $C_n$ for $n = 3,\ldots,12$, but a general proof
requires an argument independent of the size of $\Gloc$.
\end{openproblem}

\begin{openproblem}{Criterion for the Ambiguous Region}{OP-SP2-2}
Rule~A fails for materials with $\DEN < 1.0$: GaAs (direct) and
GaP (indirect) have nearly identical $\DEN$ (0.37 vs.\ 0.38) and
cannot be distinguished by topological means alone. Identify
an additional PDL-derivable criterion — possibly related to orbital
character at band edges — that resolves this ambiguity.
\end{openproblem}

\begin{openproblem}{Extension to Molecular Frameworks}{OP-SP2-3}
Apply PDL-H to covalent organic frameworks (COF) based on metalloporphyrins.
These systems have non-bipartite local chromophores (girth~$= 5$)
organised in hexagonal 2D networks (bipartite globally). PDL-H predicts
$\epsgeom(H) = 0$; this should be verified on crystallographic structures
from the Cambridge Structural Database, and compared with measured
quantum yields.
\end{openproblem}

%% ══════════════════════════════════════════════════════════════
\section{Computational Resources}
\label{sec:computational}

All numerical verifications were performed in Google Colaboratory
using Python 3.10 with NetworkX 3.x, pymatgen, and the Materials
Project API (mp-api). Crystallographic data retrieved from the
Materials Project \cite{Jain2013}. Biological structure 3ARC from
the RCSB Protein Data Bank \cite{Umena2011}, parsed with Biopython.
All computations use exact arithmetic (Python \texttt{Fraction} class)
where applicable. The following Colab notebooks constitute the
complete computational trace of this document:

\begin{enumerate}[label=\textbf{NB\arabic*.}]
  \item \texttt{PDL\_OP\_SP2\_derivation.ipynb} — Theorem PDL-SP2,
    bijection verification, counting formula, connectivity necessity.
  \item \texttt{PDL\_OP\_SP2\_proof.ipynb} — Analytic proof verification
    (Part~A telescoping, Part~B BFS reconstruction, non-connected case).
  \item \texttt{PDL\_OP\_SP2\_clean.ipynb} — Material screening (7
    candidates, covalent cutoff, bipartition, dynamic formula lookup).
  \item \texttt{PDL\_OP\_SP2\_extended.ipynb} — Extended screening
    (21 materials, 4 families, composite score 0--4).
  \item \texttt{PDL\_OP\_SP2\_stepA.ipynb} — Rule~A derivation
    ($\DEN$ vs.\ direct/indirect gap, 15 materials).
  \item \texttt{PDL\_OP\_SP2\_stepA\_corrected.ipynb} — Corrected
    supercell analysis (girth and $k$ with 2$\times$2$\times$2 cells).
  \item \texttt{PDL\_OP\_SP2\_stepB.ipynb} — Biological verification
    on Photosystem~II structure 3ARC (chlorophyll, Mg--Mg network).
  \item \texttt{PDL\_OP\_SP2\_hierarchical.ipynb} — Hierarchical graph
    analysis (Verification~1, all combinations).
  \item \texttt{PDL\_OP\_SP2\_verification2.ipynb} — Verification~2
    ($C_n$ for $n=3,\ldots,12$, pair/odd frontier).
  \item \texttt{PDL\_OP\_SP2\_exhaustive\_C5C4.ipynb} — Exhaustive
    verification $H(C_5, C_4)$, $2^{24}$ assignments, 20~minutes
    computation time, $\epsgeom = 0$ confirmed.
\end{enumerate}

All notebooks are archived with this document on Zenodo
(community: \texttt{pdl-framework}).
 
%% ── References ────────────────────────────────────────────────
\bibliographystyle{unsrt}
\bibliography{D-exp-SP2}

\end{document}
